\section{Experiments}

\subsection{Existence Gate Experiment (h-e1)}

We test the first causal mechanism step: do static analyzers produce actionable warnings on benchmark code?

\textbf{Setup:}
\begin{itemize}
    \item Dataset: HumanEval (164 problems, full test set)
    \item Code source: Canonical solutions (proxy for LLM output due to API unavailability)
    \item Analyzer: pylint 3.x with \texttt{--disable=C,R}
    \item Gate threshold: $\geq$30\% of problems with $\geq$1 warning
\end{itemize}

\textbf{Rationale for Proxy:} Without OpenAI/Claude API access, we used canonical solutions as a lower-bound proxy. Canonical solutions are well-written human code; LLM-generated code typically exhibits more issues.

\subsection{Implementation}

\textbf{Pipeline Components:}
\begin{enumerate}
    \item \texttt{humaneval\_loader.py} --- Loads HumanEval problems from HuggingFace
    \item \texttt{static\_analyzer.py} --- Runs pylint, parses JSON output, filters by severity
    \item \texttt{metrics.py} --- Aggregates warning counts, computes rates
\end{enumerate}

\textbf{Pylint Configuration:}
\begin{verbatim}
pylint --output-format=json --disable=C,R <code_file>
\end{verbatim}

\subsection{Artifacts}

\begin{table}[h]
\centering
\caption{Generated artifacts from the existence gate experiment.}
\label{tab:artifacts}
\begin{tabular}{@{}ll@{}}
\toprule
Artifact & Description \\
\midrule
\texttt{metrics.json} & Aggregated warning statistics \\
\texttt{pylint\_results.json} & Per-problem pylint output \\
\texttt{warning\_distribution.png} & Histogram of warnings per problem \\
\texttt{warning\_type\_breakdown.png} & Pie chart by category \\
\texttt{gate\_metrics.png} & Threshold vs.\ actual comparison \\
\texttt{top\_warning\_codes.png} & Top 10 warning codes \\
\bottomrule
\end{tabular}
\end{table}
